\documentclass[10pt,twocolumn,letterpaper]{article}
\usepackage{cvpr}
\usepackage{times}
\usepackage{graphicx}
\usepackage{amsmath,amssymb}
\def\paperID{*****}
\def\confName{CVPR}
\def\confYear{2027}
\begin{document}
\onecolumn
\noindent\includegraphics[width=\textwidth]{design/benchmark_overview.pdf}

{\small\noindent\textbf{Figure 1. ContextWorld pairs interaction histories to test dynamics adaptation.} A: Real PushT observations show two action gains. The current observation and future actions are identical, but the true futures differ. All images use the same crop; colored arrows mark the pusher's displacement from the shared current position. B: The same frozen model makes separate latent predictions. Matrix entries are MSEs between predicted and encoded true futures. Rows compare possible futures; columns compare histories for a fixed target. The displayed inequalities illustrate condition A; both conditions are scored.}

\vspace{14pt}
\noindent\begin{minipage}[t]{0.48\textwidth}
\vspace{0pt}
\includegraphics[width=\linewidth]{design/training_two_panels.pdf}

{\small\textbf{Figure 4. Prediction loss and conditional response.} PLDM is trained on contact friction and motion damping, with the visual encoder jointly updated or frozen. Loss is normalized to epoch 5; bands show batch interquartile ranges. Training-query NRE is measured at epochs 5 and 10 only; 1 marks the zero-response reference. Within each task, both schemes start from the same Original checkpoint.}
\end{minipage}\hfill\begin{minipage}[t]{0.48\textwidth}
\vspace{0pt}
\textbf{Prediction Loss and Adaptation}

Weak evaluation performance may reflect poor generalization or conditional responses that were not learned well during training. We compare prediction loss with NRE on fixed training queries at epochs 5 and 10.

On friction, PLDM joint training reduces prediction loss by 18.05\%, yet training-query NRE remains near the zero-response reference. On damping with the encoder frozen, loss approaches convergence while NRE improves but remains high. Changing the scheme helps: freezing yields a stronger conditional response on friction, whereas joint training performs better on damping.

These results expose a disconnect between prediction-loss optimization and history-conditioned adaptation, even on training data. Prediction loss aggregates errors across mixed data, whereas NRE directly measures error in the predicted differences between conditions.
\end{minipage}
\end{document}
